% Opening from lecture 05; extend with subsequent lectures on sensing.
\chapter{Robot systems and sensing}
\label{chap:lecture06}

Reference frames and transformations describe where a robot and the objects
around it are located. A robot must also obtain information about those objects,
decide what to do, and produce physical motion. These activities connect the
geometric tools of the preceding chapters to the operation of a robot system.

\section{Sensing, processing, and action}
A useful working description of a robot is an intelligent interconnection
between sensing and action. Sensors collect measurements; computation uses
those measurements to select actions; actuators carry out the actions.
The resulting interaction with the environment produces new measurements,
closing a feedback loop.

\begin{figure}[htbp]
\centering\begin{tikzpicture}[box/.style={draw,rounded corners,align=center,minimum height=1cm,minimum width=2.6cm}]
\draw[rounded corners,robotframe,thick] (-2.1,-2) rectangle (5.4,2);
\node[robotframe,fill=white] at (1.8,2) {Robot};
\node[box] (I) at (0,0) {Processing, planning,\\and control};
\node[box] (A) at (3.7,1) {Actuators};
\node[box] (S) at (3.7,-1) {Sensors};
\node[box,minimum width=2cm,minimum height=3cm] (E) at (7.2,0) {Environment};
\draw[axis] (I.north) |- (A.west) node[pos=.7,above] {commands};
\draw[axis] (S.west) -| (I.south) node[pos=.35,below] {measurements};
\draw[axis] (A.east)--(A.east -| E.west) node[midway,above] {action};
\draw[axis] (S.east -| E.west)--(S.east);
\end{tikzpicture}

\caption{A functional view of a robot interacting with its environment.
The robot boundary contains sensors, actuators, and the hardware and software
that process measurements and generate commands.}
\label{fig:l06-components}
\end{figure}

\section{Sensors and the environment}
The \emph{environment} includes the objects, surfaces, and people with which
the robot interacts. A camera measures incoming light, a microphone responds
to pressure variations, and a contact sensor detects physical contact.
Sensors may also measure the robot's own state; sensing is not limited to
observing external objects.

A measurement is not yet an interpretation of the scene. An image records
light arriving at the camera; identifying an object and estimating its
location require processing that image. Likewise, a contact measurement
indicates an interaction at a surface but does not by itself describe the
entire surrounding space. Each type of sensor provides particular information
that the robot must interpret in relation to its task.

The coordinate transformations developed earlier enter at this point. An
object located in a sensor frame may need to be expressed in a robot frame
to determine whether it is within reach, or in a world frame to plan motion.
The physical object is the same; the coordinate description is chosen to
support the calculation being performed.

\section{Actuators and physical interaction}
\emph{Actuators} produce physical action, for example by driving wheels or
moving a manipulator. Wheel-driven motion depends on forces exchanged with
the supporting surface. An action may change the environment, the robot's
pose relative to it, or both.

The distinction between sensing and action is a useful functional one:
sensors provide measurements, whereas actuators produce physical effects.
A motor command can turn a wheel or move a joint, but the resulting motion
also depends on the mechanism and its interaction with the surroundings.
Moving a mobile robot changes what its sensors observe, even when the objects
in the environment remain fixed in the world frame.

\section{Connecting measurements to commands}
The connection between sensing and action includes both computing hardware
and software. It may involve interpreting measurements, estimating where
objects are, selecting a plan, and producing individual actuator commands.
Here, ``intelligent'' does not require a neural network or machine learning:
explicit algorithms and feedback control can also provide this connection.

The central block in Figure~\ref{fig:l06-components} therefore represents
several possible levels of computation. A high-level decision might be to
approach an object. More detailed decisions specify how to move, and the
commands sent to the actuators must ultimately specify actions that the
hardware can execute. All of these functions belong to the connection
between measurements and physical action, even when they run in different
programs or on different processors.

The enclosing boundary in the diagram is also important: the robot includes
the sensors and actuators as well as the computing hardware and software.
The central block alone is not the entire robot. The arrows trace a loop
through the environment: an action changes the physical situation, sensors
measure the resulting situation, and computation can use those measurements
to decide the next action.

\subsection{Example: reaching an object}
Suppose a robot must grasp an object that is initially beyond its reach.
Sensor measurements first support an estimate of the object's location.
The robot must then determine that it needs to move closer, choose an approach,
and command its motion before extending the manipulator and closing the gripper.
A high-level objective thus becomes a sequence of physical actions, with
measurements available to check progress and adjust the commands.

Consider the distinction between the decisions involved. Recognizing the
object answers what the robot should interact with. Estimating its position
answers where it is. Comparing that position with the robot's reach determines
whether the base must move first. Planning selects a sequence of motions,
and control turns those motions into actuator commands. These decisions
serve the same task but require different information and computations.

After the robot approaches, the object's coordinates relative to the robot
have changed. New measurements can update its estimated location before the
grasp. This illustrates why sensing and action are connected repeatedly:
the robot does not merely observe once and then execute an unrelated motion.
